\documentclass[12pt,a4paper]{article}

% ---- Encoding & Font ----
\usepackage{iftex}
\ifPDFTeX
  \usepackage[utf8]{inputenc}
  \usepackage[T1]{fontenc}
\else
  \usepackage{fontspec}  % XeTeX/LuaTeX (tectonic): Unicode nativo (·, —, ¿, ª, §)
\fi
\usepackage[spanish]{babel}
\usepackage{csquotes}

% ---- Page layout ----
\usepackage[top=2.5cm, bottom=2.5cm, left=2.5cm, right=2.5cm]{geometry}
\usepackage{setspace}
\onehalfspacing

% ---- Math ----
\usepackage{amsmath, amssymb, amsthm}
\usepackage{mathtools}
\usepackage{physics}
\usepackage{esint}  % \oiint

% ---- Graphics ----
\usepackage{graphicx}
\usepackage{tikz}
\usepackage{float}
\usepackage{caption}

% ---- Tables ----
\usepackage{array}
\usepackage{booktabs}
\usepackage{tabularx}
\usepackage{colortbl}

% ---- Colours ----
\usepackage{xcolor}
\definecolor{notebg}{HTML}{FFF3CD}
\definecolor{noteborder}{HTML}{FFC107}
\definecolor{boxred}{HTML}{F8D7DA}
\definecolor{boxredborder}{HTML}{F5C6CB}

% ---- Boxes ----
\usepackage{mdframed}
\usepackage{tcolorbox}
\tcbuselibrary{most}

\newtcolorbox{keybox}[1][]{
  colback=blue!5,
  colframe=blue!70!black,
  arc=4pt,
  boxrule=1pt,
  fonttitle=\bfseries,
  title={#1}
}

\newtcolorbox{warnbox}[1][]{
  colback=boxred,
  colframe=boxredborder,
  coltitle=black!75,
  arc=4pt,
  boxrule=1pt,
  fonttitle=\bfseries,
  title={#1}
}

\newtcolorbox{notebox}[1][]{
  colback=notebg,
  colframe=noteborder,
  coltitle=black!75,
  arc=4pt,
  boxrule=1pt,
  fonttitle=\bfseries,
  title={#1}
}

% ---- Hyperlinks & TOC ----
\usepackage[hidelinks]{hyperref}
\usepackage{bookmark}

% ---- Enumerate ----
\usepackage{enumitem}

% ---- Miscellaneous ----
\usepackage{icomma}
\usepackage{siunitx}
\usepackage{textcomp}
\usepackage[bottom]{footmisc}

% ---- Diagramas (gráficas y esquemas) ----
\usepackage{pgfplots}
\pgfplotsset{compat=1.18}
\usetikzlibrary{babel}  % babel español activa `>`; esto evita romper las flechas `->`
% courses/ElecMag2024/Clases/assets/tikzstyles.tex
% ---------------------------------------------------------------------------
% Estilos compartidos para los diagramas del curso Electromagnetismo (2024).
% Fuente única del "look" visual (colores, grosores, escalas) para que todas
% las figuras (TikZ / pgfplots / CircuiTikz) sean consistentes entre clases.
%
% Es una COPIA de courses/Fisica3-2015/Clases/assets/tikzstyles.tex: mismo
% dominio, misma paleta. Copia y no symlink para que cada curso sea
% autocontenido y la paleta de Física III pueda evolucionar aparte.
%
% Uso: la clase debe cargar antes `tikz` y `pgfplots` (y `circuitikz` si la
%      clase tiene circuitos), y luego % courses/ElecMag2024/Clases/assets/tikzstyles.tex
% ---------------------------------------------------------------------------
% Estilos compartidos para los diagramas del curso Electromagnetismo (2024).
% Fuente única del "look" visual (colores, grosores, escalas) para que todas
% las figuras (TikZ / pgfplots / CircuiTikz) sean consistentes entre clases.
%
% Es una COPIA de courses/Fisica3-2015/Clases/assets/tikzstyles.tex: mismo
% dominio, misma paleta. Copia y no symlink para que cada curso sea
% autocontenido y la paleta de Física III pueda evolucionar aparte.
%
% Uso: la clase debe cargar antes `tikz` y `pgfplots` (y `circuitikz` si la
%      clase tiene circuitos), y luego % courses/ElecMag2024/Clases/assets/tikzstyles.tex
% ---------------------------------------------------------------------------
% Estilos compartidos para los diagramas del curso Electromagnetismo (2024).
% Fuente única del "look" visual (colores, grosores, escalas) para que todas
% las figuras (TikZ / pgfplots / CircuiTikz) sean consistentes entre clases.
%
% Es una COPIA de courses/Fisica3-2015/Clases/assets/tikzstyles.tex: mismo
% dominio, misma paleta. Copia y no symlink para que cada curso sea
% autocontenido y la paleta de Física III pueda evolucionar aparte.
%
% Uso: la clase debe cargar antes `tikz` y `pgfplots` (y `circuitikz` si la
%      clase tiene circuitos), y luego \input{../assets/tikzstyles.tex}
% (compilar desde el directorio de la clase para que la ruta relativa resuelva).
% ---------------------------------------------------------------------------

% ---- Paleta (alineada con las cajas del preámbulo: keybox/notebox/warnbox) ----
\definecolor{figblue}{HTML}{1F4E79}   % azul principal (líneas, keybox)
\definecolor{figred}{HTML}{C0392B}    % rojo de énfasis (warnbox)
\definecolor{figamber}{HTML}{B8860B}  % ámbar (notebox)
\definecolor{figgray}{HTML}{555555}   % auxiliares, ejes, guías

% ---- CircuiTikz: escala y grosor de los componentes ----
% Guarda \ifdefined: la electrostática (Clases 1-14) no carga `circuitikz`, y
% sin la guarda el \input revienta con `Undefined control sequence \ctikzset`.
% Cuando el curso llegue a circuitos, el bloque se activa solo.
\ifdefined\ctikzset
  \ctikzset{
    bipoles/thickness=1,
    resistors/scale=0.9,
    capacitors/scale=0.9,
  }
\fi

% ---- TikZ: estilos reutilizables ----
% Nota: las puntas se declaran como `-{Latex[...]}` y no como `->`. La forma
% con llaves NO contiene un `>` literal, así que es inmune al `>` activo que
% introduce babel español (ver CLAUDE.md §7.3.3).
\usetikzlibrary{arrows.meta,calc,angles,quotes}

\tikzset{
  figline/.style={figblue, line width=0.9pt},
  figaux/.style={figgray, dashed, line width=0.6pt},
  % vectores
  figvec/.style={figblue, line width=0.9pt, -{Latex[length=2.2mm,width=1.6mm]}},
  figvecres/.style={figred, line width=1.3pt, -{Latex[length=2.6mm,width=1.9mm]}},
  figvecaux/.style={figgray, line width=0.7pt, -{Latex[length=1.8mm,width=1.3mm]}},
  figaxis/.style={figgray, line width=0.6pt, -{Latex[length=2mm,width=1.4mm]}},
  % cargas puntuales
  figcarga/.style={circle, inner sep=0pt, minimum size=5.4pt, draw=none},
  figpos/.style={figcarga, fill=figred},
  figneg/.style={figcarga, fill=figblue},
  figneutra/.style={figcarga, fill=figgray},
  % etiquetas
  figlbl/.style={font=\small},
  figlblaux/.style={font=\footnotesize, figgray},
}

% ---- pgfplots: estilo base de las gráficas ----
\pgfplotsset{
  emfig/.style={
    width=11cm, height=6.5cm,
    axis lines=left,
    xlabel style={font=\small},
    ylabel style={font=\small},
    tick label style={font=\footnotesize},
    every axis plot/.append style={line width=1pt},
    grid=major,
    grid style={figgray!20},
    legend cell align=left,
    legend style={font=\footnotesize, draw=figgray!40},
  },
}

% (compilar desde el directorio de la clase para que la ruta relativa resuelva).
% ---------------------------------------------------------------------------

% ---- Paleta (alineada con las cajas del preámbulo: keybox/notebox/warnbox) ----
\definecolor{figblue}{HTML}{1F4E79}   % azul principal (líneas, keybox)
\definecolor{figred}{HTML}{C0392B}    % rojo de énfasis (warnbox)
\definecolor{figamber}{HTML}{B8860B}  % ámbar (notebox)
\definecolor{figgray}{HTML}{555555}   % auxiliares, ejes, guías

% ---- CircuiTikz: escala y grosor de los componentes ----
% Guarda \ifdefined: la electrostática (Clases 1-14) no carga `circuitikz`, y
% sin la guarda el \input revienta con `Undefined control sequence \ctikzset`.
% Cuando el curso llegue a circuitos, el bloque se activa solo.
\ifdefined\ctikzset
  \ctikzset{
    bipoles/thickness=1,
    resistors/scale=0.9,
    capacitors/scale=0.9,
  }
\fi

% ---- TikZ: estilos reutilizables ----
% Nota: las puntas se declaran como `-{Latex[...]}` y no como `->`. La forma
% con llaves NO contiene un `>` literal, así que es inmune al `>` activo que
% introduce babel español (ver CLAUDE.md §7.3.3).
\usetikzlibrary{arrows.meta,calc,angles,quotes}

\tikzset{
  figline/.style={figblue, line width=0.9pt},
  figaux/.style={figgray, dashed, line width=0.6pt},
  % vectores
  figvec/.style={figblue, line width=0.9pt, -{Latex[length=2.2mm,width=1.6mm]}},
  figvecres/.style={figred, line width=1.3pt, -{Latex[length=2.6mm,width=1.9mm]}},
  figvecaux/.style={figgray, line width=0.7pt, -{Latex[length=1.8mm,width=1.3mm]}},
  figaxis/.style={figgray, line width=0.6pt, -{Latex[length=2mm,width=1.4mm]}},
  % cargas puntuales
  figcarga/.style={circle, inner sep=0pt, minimum size=5.4pt, draw=none},
  figpos/.style={figcarga, fill=figred},
  figneg/.style={figcarga, fill=figblue},
  figneutra/.style={figcarga, fill=figgray},
  % etiquetas
  figlbl/.style={font=\small},
  figlblaux/.style={font=\footnotesize, figgray},
}

% ---- pgfplots: estilo base de las gráficas ----
\pgfplotsset{
  emfig/.style={
    width=11cm, height=6.5cm,
    axis lines=left,
    xlabel style={font=\small},
    ylabel style={font=\small},
    tick label style={font=\footnotesize},
    every axis plot/.append style={line width=1pt},
    grid=major,
    grid style={figgray!20},
    legend cell align=left,
    legend style={font=\footnotesize, draw=figgray!40},
  },
}

% (compilar desde el directorio de la clase para que la ruta relativa resuelva).
% ---------------------------------------------------------------------------

% ---- Paleta (alineada con las cajas del preámbulo: keybox/notebox/warnbox) ----
\definecolor{figblue}{HTML}{1F4E79}   % azul principal (líneas, keybox)
\definecolor{figred}{HTML}{C0392B}    % rojo de énfasis (warnbox)
\definecolor{figamber}{HTML}{B8860B}  % ámbar (notebox)
\definecolor{figgray}{HTML}{555555}   % auxiliares, ejes, guías

% ---- CircuiTikz: escala y grosor de los componentes ----
% Guarda \ifdefined: la electrostática (Clases 1-14) no carga `circuitikz`, y
% sin la guarda el \input revienta con `Undefined control sequence \ctikzset`.
% Cuando el curso llegue a circuitos, el bloque se activa solo.
\ifdefined\ctikzset
  \ctikzset{
    bipoles/thickness=1,
    resistors/scale=0.9,
    capacitors/scale=0.9,
  }
\fi

% ---- TikZ: estilos reutilizables ----
% Nota: las puntas se declaran como `-{Latex[...]}` y no como `->`. La forma
% con llaves NO contiene un `>` literal, así que es inmune al `>` activo que
% introduce babel español (ver CLAUDE.md §7.3.3).
\usetikzlibrary{arrows.meta,calc,angles,quotes}

\tikzset{
  figline/.style={figblue, line width=0.9pt},
  figaux/.style={figgray, dashed, line width=0.6pt},
  % vectores
  figvec/.style={figblue, line width=0.9pt, -{Latex[length=2.2mm,width=1.6mm]}},
  figvecres/.style={figred, line width=1.3pt, -{Latex[length=2.6mm,width=1.9mm]}},
  figvecaux/.style={figgray, line width=0.7pt, -{Latex[length=1.8mm,width=1.3mm]}},
  figaxis/.style={figgray, line width=0.6pt, -{Latex[length=2mm,width=1.4mm]}},
  % cargas puntuales
  figcarga/.style={circle, inner sep=0pt, minimum size=5.4pt, draw=none},
  figpos/.style={figcarga, fill=figred},
  figneg/.style={figcarga, fill=figblue},
  figneutra/.style={figcarga, fill=figgray},
  % etiquetas
  figlbl/.style={font=\small},
  figlblaux/.style={font=\footnotesize, figgray},
}

% ---- pgfplots: estilo base de las gráficas ----
\pgfplotsset{
  emfig/.style={
    width=11cm, height=6.5cm,
    axis lines=left,
    xlabel style={font=\small},
    ylabel style={font=\small},
    tick label style={font=\footnotesize},
    every axis plot/.append style={line width=1pt},
    grid=major,
    grid style={figgray!20},
    legend cell align=left,
    legend style={font=\footnotesize, draw=figgray!40},
  },
}


% ---- Title ----
\title{\textbf{Clase 12: Energía electrostática} \\[0.3em]
        \large{Distribuciones continuas, densidad de energía y conductores cargados}}
\author{Transcripto y expandido --- Curso Electromagnetismo (OpenFING)}
\date{}

\begin{document}

\maketitle
\thispagestyle{empty}
\newpage

\tableofcontents
\newpage

\section{Repaso: energía de cargas puntuales}

\subsection{Parejas y pares}

La clase anterior cerró con la energía de un sistema de $N$ cargas puntuales,
donde $r_{ij}$ es la distancia entre las cargas $i$ y $j$ (por ejemplo, $r_{14}$
entre la 1 y la 4):
\[
U = \frac{1}{4\pi\varepsilon_0}\sum_{i=1}^{N}\sum_{j=1}^{i-1}\frac{q_i q_j}{r_{ij}} .
\qquad(\ast\ast)
\]
Es un resultado de Física 3. Se suma sobre todas las \textbf{parejas}, cada una
una sola vez: con tres cargas son tres términos, $(1,2)$, $(1,3)$ y $(2,3)$.

\begin{figure}[H]
  \centering
  % Clase 12 — cuatro cargas puntuales y las distancias entre ellas.
% Se dibujan las SEIS parejas (4 sobre 2) porque ese es el contenido: la suma va
% sobre parejas, no sobre "una con todas". Solo dos distancias rotuladas; las
% seis juntas saturaban el dibujo y no agregaban nada.
\begin{tikzpicture}[scale=1]

  \coordinate (Q1) at (0,0);
  \coordinate (Q2) at (3.2,1.9);
  \coordinate (Q3) at (5.4,-0.3);
  \coordinate (Q4) at (2.0,-2.0);

  % parejas
  \foreach \a/\b in {Q1/Q2,Q1/Q3,Q2/Q3,Q2/Q4,Q3/Q4}{
    \draw[figaux] (\a) -- (\b);}
  \draw[figline] (Q1) -- (Q4) node[midway, left, figlbl, figblue, xshift=-2pt] {$r_{14}$};
  \draw[figline] (Q2) -- (Q3) node[midway, right, figlbl, figblue, xshift=3pt, yshift=3pt] {$r_{23}$};

  % cargas
  \node[figpos] at (Q1) {}; \node[figlbl] at ($(Q1)+(-0.35,0.25)$) {$q_1$};
  \node[figpos] at (Q2) {}; \node[figlbl] at ($(Q2)+(0,0.35)$) {$q_2$};
  \node[figneg] at (Q3) {}; \node[figlbl] at ($(Q3)+(0.4,0.1)$) {$q_3$};
  \node[figpos] at (Q4) {}; \node[figlbl] at ($(Q4)+(0,-0.4)$) {$q_4$};

\end{tikzpicture}

  \caption{Cuatro cargas, seis parejas. Cada pareja aporta un término
  $q_iq_j/r_{ij}$ a la energía, y se cuenta una sola vez.}
\end{figure}

La misma cantidad se puede escribir sumando sobre todos los \textbf{pares
ordenados} ---contando $(1,2)$ y $(2,1)$ por separado--- y dividiendo por dos
para no contar doble:
\[
U = \frac{1}{2}\,\frac{1}{4\pi\varepsilon_0}\sum_{i=1}^{N}\sum_{\substack{j=1\\ j\neq i}}^{N}\frac{q_i q_j}{r_{ij}} .
\]

\begin{warnbox}[No se cuenta la autointeracción]
El término $i=j$ daría infinito, porque la distancia de una carga a sí misma es
cero. La energía es \textbf{interacción entre cosas}, no de una cosa consigo
misma. El docente cita un dicho que nació precisamente de este problema:
\emph{no porque algo sea infinito es forzosamente cero}. Sobre esto se vuelve en
§\ref{sec:autoenergia}.
\end{warnbox}

\subsection{La energía en términos de los potenciales}

El potencial en la posición de la carga $i$ es el que generan \textbf{todas las
demás}:
\[
\phi_i = \frac{1}{4\pi\varepsilon_0}\sum_{j\neq i}\frac{q_j}{r_{ij}} .
\]
Separando la suma en $i$ en la fórmula de pares, lo que queda es justamente
$\phi_i$:

\begin{keybox}[Energía en términos de cargas y potenciales]
\[
\boxed{\ U = \frac{1}{2}\sum_{i=1}^{N} q_i\,\phi_i\ }\qquad(\ast)
\]
\end{keybox}

\begin{table}[H]
  \centering
  \begin{tabularx}{\textwidth}{lX}
    \toprule
    Fórmula & Vale en \\
    \midrule
    $(\ast\ast)$ --- en términos de $q_iq_j/r_{ij}$ & sólo en el \textbf{vacío} \\
    $(\ast)$ --- en términos de $q_i\phi_i$ & también en presencia de dieléctricos \textbf{homogéneos, isótropos y lineales} \\
    \bottomrule
  \end{tabularx}
\end{table}

\begin{notebox}[Todavía no está probado]
La deducción anterior de $(\ast)$ sólo vale en el vacío, porque parte de
$(\ast\ast)$. Que $(\ast)$ sea más general se prueba ahora, para distribuciones
continuas. Para cargas puntuales la prueba es idéntica y queda como ejercicio.
\end{notebox}

\section{Energía de una distribución continua}

Sección 4.2. Se considera una densidad volumétrica $\rho$ y una superficial
$\sigma$ (podrían agregarse densidades lineales o cargas puntuales; la prueba es
la misma). Cuando no se aclara, las densidades son las \textbf{libres}
(externas). \textbf{Hipótesis:} dieléctricos lineales.

\subsection{El proceso de carga}

La energía es el \textbf{trabajo necesario para instaurar la distribución desde
cero}. Se toma $\phi=0$ cuando $\rho\equiv 0$ y $\sigma\equiv 0$, y se van
incrementando las cargas \textbf{cuasiestáticamente} ---«lento, lento, lento»---
hasta sus valores finales.

\begin{notebox}[«Es electrostática, nada se mueve»]
Sí, pero si nada se moviera nunca se llegaría a la configuración deseada. Se la
arma de a poquito, y se calcula la energía que cuesta ese proceso.
\end{notebox}

En un estado intermedio con carga $q'$ y potencial $\phi'$, agregar $\delta q$
cuesta
\[
\delta W = \phi'(\vec r)\,\delta q,
\qquad
\delta q = \delta\rho\,dV \ \ \text{(volumen)}
\quad\text{o}\quad
\delta q = \delta\sigma\,dS \ \ \text{(superficie)}.
\]
Como el sistema es \textbf{conservativo}, el resultado no depende del camino por
el que se suban las cargas: se podría subir primero la volumétrica y después la
superficial, o unas más rápido que otras. Se elige el camino más cómodo:
\textbf{todas crecen proporcionalmente},
\[
\rho' = \alpha\,\rho,\qquad \sigma' = \alpha\,\sigma,\qquad \alpha\in[0,1],
\]
con $\alpha=0$ al principio y $\alpha=1$ en la configuración final. Entonces
$\delta\rho = \rho\,\delta\alpha$ y $\delta\sigma = \sigma\,\delta\alpha$, y
\[
U = \int_0^1 d\alpha\left[\int_V \phi'(\vec r)\,\rho(\vec r)\,dV
+ \int_S \phi'(\vec r)\,\sigma(\vec r)\,dS\right].
\]

\begin{notebox}[Qué queda adentro y qué afuera]
Son integrales una dentro de la otra. $\rho$ y $\sigma$ quedan \textbf{adentro}
porque dependen de $\vec r$; $d\alpha$ sale \textbf{afuera} porque se eligió el
mismo para todos los puntos. Si no se hubiera elegido que todas crecen igual, eso
no se podría hacer.
\end{notebox}

\subsection{Linealidad y el factor un medio}

Acá entra la linealidad del medio: el potencial es \textbf{proporcional a la
carga}. Si todas las cargas se multiplican por un factor, el potencial se
multiplica por el mismo factor:
\[
\phi'(\vec r) = \alpha\,\phi(\vec r),
\]
con $\phi$ el potencial final. Toda la dependencia en $\alpha$ queda explícita, y
$\int_0^1 \alpha\,d\alpha = 1/2$:

\begin{keybox}[Energía de una distribución continua]
\[
\boxed{\ U = \frac12\int_V \rho\,\phi\,dV + \frac12\int_S \sigma\,\phi\,dS\ }
\]
Es la contraparte continua de $U = \tfrac12\sum_i q_i\phi_i$.
\end{keybox}

\begin{notebox}[Lo que se usó, y lo que no]
Sólo que el trabajo es conservativo y que el medio es lineal. Es una prueba
completamente distinta de la de cargas puntuales, pero da el mismo resultado; el
vacío es un caso particular de medio lineal. Y el $\tfrac12$ viene \textbf{sólo}
de la integral en $\alpha$.
\end{notebox}

\subsection{La paradoja de la autoenergía}
\label{sec:autoenergia}

Comparar esta fórmula con $(\ast)$ debería generar «cierta incomodidad». En
$(\ast)$, $\phi_i$ es el potencial de \textbf{las otras} cargas; la propia no
cuenta. En la fórmula continua no se hizo esa distinción: $\phi$ es el potencial
total.

Una carga puntual es una \textbf{idealización} de una distribución muy
concentrada, así que la fórmula de cargas puntuales debería salir de la continua
en el límite. No sale: se obtiene $(\ast)$ \textbf{más} la energía de interacción
interna de cada carga consigo misma. Si se imagina la carga como una esferita, esa
energía no tiende a cero al achicarla sino que \textbf{diverge}: el volumen baja,
pero el $1/r$ crece.

\begin{keybox}[Resolución]
Lo que interesa es la energía de \textbf{interacción} entre partículas, que
depende de sus \textbf{distancias relativas}. La energía interna de cada carga no
depende de su posición: es una constante aditiva ---la suma de las
\textbf{autoenergías}---. Como la energía potencial está definida a menos de una
constante, se la descarta cambiando el nivel cero. Y es bueno descartarla, porque
en el límite puntual es infinita.
\end{keybox}

\begin{notebox}[Nota de cultura general]
Un problema análogo, pero en mecánica cuántica, ocupó a los mejores físicos del
mundo aproximadamente entre 1924 y 1947 y su resolución dio lugar a un premio
Nobel. En el caso clásico es trivial ---es sólo una constante---; en el cuántico
las autoenergías se mezclan con la dinámica. El concepto de la solución es el
mismo: en las cantidades realmente observables, como distancias relativas, esos
infinitos no aparecen. El electrón, en los experimentos más precisos, no muestra
estructura interna, y con la carga cuantizada no queda claro qué estructura podría
tener algo cuya carga ya es la mínima posible. Para este curso, nada de esto
importa.
\end{notebox}

\subsection{Conductores}
\label{sec:conductores}

En un conductor $\rho = 0$ y sólo hay $\sigma$ en su superficie, donde el
potencial es \textbf{constante} en situación estática. Sale de la integral:
\[
U = \frac12\,\phi_c\oint_S\sigma\,dS
\qquad\Longrightarrow\qquad
\boxed{\ U = \frac12\,Q\,\phi_c\ }
\]

\begin{notebox}
Una fórmula como la de una carga puntual ---potencial por carga--- aunque la
carga esté distribuida. Y \textbf{no depende de cómo está distribuida}, sólo de
la total, lo cual simplifica mucho porque muchas veces esa distribución no se
conoce.
\end{notebox}

\section{Densidad de energía electrostática}

En Física 3 se enseñó que la energía puede pensarse de dos maneras: como
interacción directa entre cargas, o como \textbf{almacenada en el campo}, con
densidad $\tfrac12\varepsilon_0 E^2$. Según el docente, aquello era «un cuentito
al borde de lo delictivo», algo que había que creer con los ojos cerrados.
Sección 4.3: se deduce ahora de manera prolija.

\subsection{Hipótesis y geometría}

\begin{enumerate}
  \item Las cargas están en una \textbf{región acotada} (no hay distribuciones que
    lleguen al infinito).
  \item Todos los dieléctricos son \textbf{lineales}.
  \item Toda la densidad superficial de carga está en \textbf{superficies de
    conductores}; no hay carga superficial sobre los dieléctricos.
\end{enumerate}

Se toma un conjunto de conductores con superficies $S_1, S_2, \dots, S_n$; entre
ellos puede haber vacío o dieléctricos lineales. Se elige además una superficie
$S'$ que encierra todo, y se llama $V$ al volumen \textbf{entre} los conductores y
$S'$. Se usan dos hechos:
\[
\div\vec D = \rho \quad\text{(Gauss)},
\qquad
\sigma = \vec D\cdot\hat n \quad\text{sobre cada } S_i ,
\]
con $\hat n$ la normal \textbf{saliente del conductor}: si $\sigma>0$, $\vec D$
sale.

\begin{figure}[H]
  \centering
  % Clase 12 — geometria de la deduccion de la densidad de energia.
% El punto delicado son las DOS normales: n sale del conductor (la que da
% sigma = D.n) y n' sale del volumen V (la del teorema de la divergencia). Sobre
% cada S_i son opuestas; sobre S' solo existe n'. Van en conductores DISTINTOS
% (n en S_1, n' en S_3): en un mismo punto, sobre S_2, se pisaban con el
% rotulo del conductor y se leian como una sola flecha. n' se dibuja desde V
% hasta la superficie, para que se vea que sale de V y entra al conductor.
\begin{tikzpicture}[scale=1]

  % S' y el volumen V
  \fill[figamber!8] (0,0) circle (3.4);
  \draw[figline, dashed] (0,0) circle (3.4);
  \node[figlbl, figblue, fill=figamber!8, inner sep=1pt] at (128:3.05) {$S'$};

  % conductores
  \fill[figgray!25, draw=figblue, line width=0.9pt]
    plot[smooth cycle, tension=0.8] coordinates {(-2.2,0.4) (-1.5,1.3) (-0.6,0.9) (-0.9,-0.1) (-1.8,-0.4)};
  \node[figlbl] at (-1.45,0.45) {$S_1$};
  \fill[figgray!25, draw=figblue, line width=0.9pt]
    plot[smooth cycle, tension=0.8] coordinates {(0.7,1.1) (1.6,1.9) (2.3,1.2) (1.6,0.5)};
  \node[figlbl] at (1.55,1.2) {$S_2$};
  \fill[figgray!25, draw=figblue, line width=0.9pt]
    plot[smooth cycle, tension=0.8] coordinates {(-0.4,-1.3) (0.6,-0.9) (1.5,-1.5) (0.6,-2.3) (-0.3,-2.1)};
  \node[figlbl] at (0.35,-1.6) {$S_3$};

  % V
  \node[figlbl, figamber!80!black] at (-1.9,-1.9) {$V$};

  % normal saliente del conductor (S_1, hacia la izquierda)
  \draw[figvecres] (-2.2,0.4) -- (-3.0,0.55);
  \node[figlbl, figred] at (-2.75,0.9) {$\hat n$};

  % normal saliente de V sobre un conductor (S_3): apunta HACIA el conductor
  \draw[figvec] (2.35,-1.55) -- (1.52,-1.5);
  \node[figlbl, figblue, fill=figamber!8, inner sep=1pt] at (1.95,-1.15) {$\hat n'$};

  % normal saliente de V sobre S'
  \draw[figvec] (-35:3.4) -- (-35:4.2);
  \node[figlbl, figblue] at (-35:4.5) {$\hat n'$};

\end{tikzpicture}

  \caption{$V$ es el volumen entre los conductores y la superficie auxiliar $S'$.
  Hay dos normales en juego: $\hat n$ sale del conductor (es la de
  $\sigma = \vec D\cdot\hat n$) y $\hat n'$ sale de $V$ (es la del teorema de la
  divergencia). Sobre cada conductor $\hat n' = -\hat n$; sobre $S'$ sólo existe
  $\hat n'$.}
\end{figure}

\subsection{El viejo truco de sumar y restar}

Sustituyendo en la energía de la sección anterior:
\[
U = \frac12\int_V (\div\vec D)\,\phi\,dV
+ \frac12\sum_{i=1}^{n}\oint_{S_i}\phi\,\vec D\cdot\hat n\,dS .
\]
Para simplificar hay que usar el teorema de la divergencia, pero el integrando no
es una divergencia: es una divergencia \textbf{por} otra cosa. El truco ---«el
viejo truco de sumar y restar», como decía el Superagente 86--- es la identidad
\[
\div(\phi\vec D) = \phi\,\div\vec D + \grad\phi\cdot\vec D .
\]
Entonces
\[
U = \frac12\int_V \div(\phi\vec D)\,dV
- \frac12\int_V \grad\phi\cdot\vec D\,dV
+ \frac12\sum_{i=1}^{n}\oint_{S_i}\phi\,\vec D\cdot\hat n\,dS .
\]
El teorema de la divergencia convierte el primer término en un flujo
\textbf{saliente de $V$} a través de \textbf{todo} su borde: las $S_i$
\textbf{y} $S'$. Con la normal $\hat n'$ saliente de $V$ ---que sobre cada
conductor es $-\hat n$--- y con $-\grad\phi = \vec E$:
\begin{multline*}
U = \frac12\int_V \vec E\cdot\vec D\,dV
+ \frac12\sum_{i=1}^{n}\oint_{S_i}\phi\,\vec D\cdot\hat n'\,dS\\
+ \frac12\oint_{S'}\phi\,\vec D\cdot\hat n'\,dS
+ \frac12\sum_{i=1}^{n}\oint_{S_i}\phi\,\vec D\cdot\hat n\,dS .
\end{multline*}
Los términos sobre los conductores se cancelan de a pares, y sobrevive sólo el de
$S'$:
\[
U = \frac12\int_V \vec E\cdot\vec D\,dV + \frac12\oint_{S'}\phi\,\vec D\cdot\hat n'\,dS .
\]

\begin{notebox}[$S'$ es un invento]
$S'$ no tiene análogo de $\hat n$: no hay conductor ahí, ni carga. Es una
superficie \textbf{inventada} para decir en qué región se calcula la energía, y
obviamente hay energía también afuera de ella.
\end{notebox}

\subsection{El término de superficie en el infinito}

Para incluir toda la energía se hace crecer $S'$: se la toma como una esfera de
radio $R$ y se manda $R\to\infty$. Como las cargas están localizadas, sobre $S'$:

\begin{table}[H]
  \centering
  \begin{tabular}{lc}
    \toprule
    Cantidad & Orden \\
    \midrule
    $\phi$ & a lo sumo $1/R$ \\
    $\lvert\vec D\rvert$ & a lo sumo $1/R^2$ \\
    área de $S'$ & $R^2$ \\
    integral sobre $S'$ & a lo sumo $1/R$ \\
    \bottomrule
  \end{tabular}
\end{table}

(Puede decrecer más rápido: si la carga total es cero, lo hace.) El término de
superficie \textbf{tiende a cero}.

\subsection{Resultado y cargas puntuales}

Queda la integral en todo el espacio fuera de los conductores. Dentro de un
conductor $\vec E = 0$, así que agregar esos volúmenes es sumar cero:

\begin{keybox}[Densidad de energía electrostática]
\[
\boxed{\ U = \frac12\int_{\text{todo el espacio}}\vec D\cdot\vec E\,dV\ },
\qquad
\boxed{\ u = \frac12\,\vec D\cdot\vec E\ }
\]
En el vacío, $\vec D = \varepsilon_0\vec E$ y se recupera la fórmula de Física 3,
$u = \tfrac12\varepsilon_0E^2$. Pero ésta es más general: vale en presencia de
dieléctricos, mientras sean lineales.
\end{keybox}

\begin{warnbox}[Ojo con las cargas puntuales]
Calculada así, la energía de un sistema con cargas puntuales da
\textbf{infinito}. Pensando la carga como una esferita de radio $r$:
$E^2\sim 1/r^4$, el volumen $\sim r^3$, y la energía diverge como $1/r$. Para
pasar a cargas puntuales hay que \textbf{sustraer la autoenergía} de cada una. La
ventaja es que no depende de las posiciones, así que no importa.
\end{warnbox}

\begin{notebox}[Tres formas, una sola energía]
Como interacción entre cargas (clase pasada); como carga por potencial
($\tfrac12\sum q_i\phi_i$ o $\tfrac12\int\rho\phi$); y como integral de una
densidad en términos de los campos. No son energías distintas: es siempre la
misma.
\end{notebox}

\section{Energía de un conjunto de conductores cargados}

Sección 4.4. El caso conocido de Física 3 es el condensador, con energía
$Q^2/2C$. Se ve ahora el caso general con $n$ conductores.

\subsection{Coeficientes de potencial}

Conductores $1,\dots,n$ con cargas $Q_1,\dots,Q_n$ y potenciales
$\phi_1,\dots,\phi_n$ en sus superficies. \textbf{Hipótesis:} todos los
dieléctricos son lineales. Se quiere probar que

\begin{keybox}[Coeficientes de potencial]
\[
\boxed{\ \phi_i = \sum_{j=1}^{n} P_{ij}\,Q_j\ }
\]
Los $P_{ij}$ dependen de la \textbf{geometría} y de las \textbf{constantes
dieléctricas}, pero \textbf{no de las cargas}: si dependieran, la relación ya no
sería lineal.
\end{keybox}

\textbf{Prueba.} Sean $\bar Q_1,\dots,\bar Q_n$ las cargas del problema que
interesa. En vez de ese problema se considera uno más simple:
\begin{enumerate}
  \item \textbf{Sólo el conductor 1 cargado}: $Q_1 = \bar Q_1$,
    $Q_2=\dots=Q_n=0$. Llamando $\phi_i^{(1)}$ al potencial resultante en el
    conductor $i$, si la carga se multiplica por $\lambda$ el potencial se
    multiplica por $\lambda$: vale Laplace, valen las condiciones de borde, y todo
    escala. Entonces $\phi_i^{(1)}$ es proporcional a $\bar Q_1$.
  \item De la misma forma se define $\phi_i^{(j)}$: el potencial en $i$ cuando
    \textbf{sólo la carga $j$} está «prendida». Es proporcional a $Q_j$.
  \item \textbf{Superposición:} con $Q_1 = \lambda\bar Q_1$, $Q_2 = \mu\bar Q_2$
    y el resto nulo, el potencial es $\lambda\phi_i^{(1)} + \mu\phi_i^{(2)}$,
    porque eso cumple las condiciones de borde. Con todas prendidas pasa lo mismo:
    \[
    \phi_i = \sum_j \phi_i^{(j)} = \sum_j P_{ij}\,Q_j ,
    \]
    donde la última igualdad es la proporcionalidad de cada $\phi_i^{(j)}$ con
    $Q_j$.
\end{enumerate}

\begin{figure}[H]
  \centering
  % Clase 12 — el paso clave de la prueba de phi_i = sum_j P_ij Q_j:
% "prender" solo una carga. El conductor 1 cargado, los demas neutros; los
% potenciales que aparecen se rotulan con el superindice de la carga prendida.
\begin{tikzpicture}[scale=1]

  % conductor 1 (cargado)
  \fill[figred!10, draw=figblue, line width=0.9pt]
    plot[smooth cycle, tension=0.8] coordinates {(-0.9,0.2) (-0.2,1.0) (0.8,0.6) (0.6,-0.5) (-0.5,-0.6)};
  \foreach \p in {(-0.75,0.25),(-0.2,0.85),(0.6,0.55),(0.5,-0.35),(-0.4,-0.45)}{
    \node[figlbl, figred] at \p {$+$};}
  \node[figlbl] at (0,0.15) {$1$};
  \node[figlbl] at (0,-1.15) {$Q_1=\bar Q_1$};
  \node[figlblaux] at (0,-1.6) {$\phi_1^{(1)}$};

  % conductor 2 (neutro)
  \fill[figgray!25, draw=figblue, line width=0.9pt]
    plot[smooth cycle, tension=0.8] coordinates {(3.1,1.4) (3.9,2.0) (4.6,1.4) (4.0,0.8)};
  \node[figlbl] at (3.9,1.4) {$2$};
  \node[figlbl] at (5.5,1.75) {$Q_2=0$};
  \node[figlblaux] at (5.5,1.3) {$\phi_2^{(1)}$};

  % conductor 3 (neutro)
  \fill[figgray!25, draw=figblue, line width=0.9pt]
    plot[smooth cycle, tension=0.8] coordinates {(3.0,-1.2) (3.8,-0.7) (4.7,-1.1) (4.3,-1.9) (3.3,-1.9)};
  \node[figlbl] at (3.85,-1.3) {$3$};
  \node[figlbl] at (5.7,-1.05) {$Q_3=0$};
  \node[figlblaux] at (5.7,-1.5) {$\phi_3^{(1)}$};

\end{tikzpicture}

  \caption{El paso 1 de la prueba: sólo el conductor 1 tiene carga. Los
  potenciales que aparecen en los tres conductores son proporcionales a
  $\bar Q_1$, y los cocientes son $P_{11}$, $P_{21}$ y $P_{31}$.}
\end{figure}

\subsection{La energía como forma cuadrática}

En un conjunto de conductores la energía es la suma de las contribuciones de cada
superficie, y en cada una el potencial sale afuera
(§\ref{sec:conductores}):
\[
U = \frac12\sum_i\oint_{S_i}\phi_i\,\sigma\,dS = \frac12\sum_i\phi_i\,Q_i .
\]
Usando la relación lineal:

\begin{keybox}[Energía de un conjunto de conductores]
\[
\boxed{\ U = \frac12\sum_{i=1}^{n}\sum_{j=1}^{n}P_{ij}\,Q_i\,Q_j\ }
\]
Aparecen todos los productos cuadráticos de todas las cargas, «de todos con
todos», con coeficientes.
\end{keybox}

\begin{notebox}[Sí aparece $i=j$]
Un conductor genera potencial sobre sí mismo: un solo conductor cargado tiene
energía. A diferencia de las cargas puntuales, acá eso no es una autoenergía
infinita.
\end{notebox}

\begin{notebox}[Cómo se calculan los $P_{ij}$]
Se apagan todas las cargas, se pone carga sólo en uno, se \textbf{resuelve
Laplace} con el método que se tenga y se miran los potenciales en todos. Con
conductores de formas raras es un problema difícil. El docente anuncia que antes
del parcial va a enseñar un método para resolver Laplace «a prueba de balas», que
siempre funciona pero necesita una computadora.
\end{notebox}

\subsection{Ejemplo: el condensador}

El caso más simple: dos placas conductoras con cargas $Q$ y $-Q$. De la fórmula
general:
\[
U = \frac12 P_{11}Q^2 + \frac12 P_{22}Q^2 - \frac12 P_{12}Q^2 - \frac12 P_{21}Q^2
= \frac12\left(P_{11}+P_{22}-P_{12}-P_{21}\right)Q^2 .
\]
Todo ese «bicho» es lo que en Física 3 se llamaba $1/C$:

\begin{keybox}[Capacitancia en términos de los coeficientes de potencial]
\[
\boxed{\ C = \frac{1}{P_{11}+P_{22}-P_{12}-P_{21}}\ },
\qquad U = \frac{Q^2}{2C}.
\]
\end{keybox}

\begin{notebox}
La energía de un condensador proporcional a $Q^2$ es un caso particular de lo
general. Y calcular la capacitancia de un condensador con placas de «formas
podridas» es el mismo problema: resolver Laplace «y remangarse».
\end{notebox}

\medskip
\noindent\emph{La clase siguiente sigue con las propiedades de los coeficientes de
potencial y una definición más general de condensador.}

\end{document}
